\begin{proof}[Beweis der endlichen Schnittstruktur\label{jordan:proof:endliche_schnittstruktur}]

    Gegeben sei eine reell-analytische Kurve $\gamma$ in der Ebene $\mathbb{R}^2$ und ein Strahl $L$ mit Anfangspunkt $p \in \mathbb{R}^2$. 
    Da $\gamma$ reell-analytisch ist, existiert eine reell-analytische Funktion $f: U \subset \mathbb{R}^2 \to \mathbb{R}$, die $\gamma$ lokal beschreibt. 
    Die Schnittpunkte von $\gamma$ und $L$ entsprechen den Nullstellen von $f$ auf $L$. 
    Da $f$ nicht verschwindend ist und $L$ eine zusammenhängende Menge ist, folgt aus dem Identitätssatz für reell-analytische Funktionen, dass die Nullstellen von $f$ auf $L$ diskret sind. 

    Entsprechend ist die Menge der Schnittpunkte $\gamma \cap L$ endlich, da eine diskrete Menge auf einer kompakten Teilmenge von $L$ nur endlich viele Punkte enthalten kann.
\end{proof}